\chapter{Kết luận và Hướng Phát triển}

\section{Tóm tắt Đóng góp}

Khóa luận đã huấn luyện và đánh giá bốn kiến trúc học sâu --- EfficientNet-B0, ResNet50, DenseNet121 và Swin-Tiny --- trên tập ISIC 2018 bảy lớp, với pipeline huấn luyện xây dựng xung quanh tình trạng mất cân bằng 58:1. Bốn thành phần giải quyết mất cân bằng từ các góc độ khác nhau: Lấy mẫu có Trọng số ở mức tải dữ liệu, Mixup và CutMix ở mức tín hiệu huấn luyện, làm mượt nhãn ở mức mục tiêu, và AdamW giảm cosine ở mức tối ưu hóa. Cùng nhau chúng loại bỏ hiện tượng sụp đổ về lớp đa số.

Ở hạt giống đại diện, EfficientNet-B0 (5,3M tham số) vượt ResNet50 (25,6M), DenseNet121 (8M) và Swin-Tiny (28M) trên mọi chỉ số phân loại cứng. Tuy nhiên khi lặp lại thí nghiệm với năm hạt giống ngẫu nhiên, Swin-Tiny mới là mô hình đơn mạnh nhất tính trung bình ($88{,}33 \pm 1{,}96\%$ accuracy so với $87{,}31 \pm 1{,}32\%$ của EfficientNet-B0); xếp hạng ở hạt giống 42 đảo ngược do phương sai của Swin cao hơn. Bài học vẫn giữ: khi dữ liệu hạn chế, hiệu quả kiến trúc và phương pháp huấn luyện quan trọng hơn số tham số, nhưng so sánh một hạt giống không đủ để xếp hạng các kiến trúc gần ngang nhau. Ensemble biểu quyết mềm nâng Macro AUC từ 0,952 lên 0,980 và F1 melanoma từ 0,695 lên 0,744 ở hạt giống đại diện, đồng thời giữ accuracy $89{,}93 \pm 0{,}66\%$ qua năm hạt giống; kiểm định McNemar xác nhận ưu thế ensemble so với mọi mô hình đơn đều có ý nghĩa thống kê ($p < 10^{-7}$ so với mô hình mạnh nhất).

Grad-CAM xác nhận mọi mô hình chú ý đến tổn thương, không phải nhiễu nền. Nguyên mẫu Streamlit tương tác chứng minh các mô hình có thể tích hợp vào giao diện hỗ trợ chẩn đoán chức năng.

Ngoài baseline chỉ-ảnh, hai thí nghiệm mở rộng được thực hiện. Kết hợp metadata bệnh nhân HAM10000 (tuổi, giới, vị trí giải phẫu) qua MLP encoder đẩy ensemble bốn mô hình lên 90,49\% accuracy và 0,873 Macro F1, vượt Pacheco và Krohling (89,0\% / 0,83) và vượt mốc 90\% --- mốc tham chiếu thường được trích dẫn cho phân loại da dựa trên ảnh. Để giới hạn kỳ vọng một cách trung thực, ensemble chỉ-ảnh cũng được đánh giá zero-shot trên PAD-UFES-20 --- 2.298 ảnh smartphone lâm sàng từ modality thu thập khác. Accuracy tụt còn 32,5\%, giảm 56 điểm phần trăm --- nhất quán với báo cáo literature cho transfer dermoscopy-sang-smartphone không adapt. AUC theo lớp vẫn giữ 0,61--0,81, xác nhận mô hình không random trên ảnh smartphone nhưng hiệu năng quyết định cứng cần adapt theo modality hoặc fine-tune để dùng được lâm sàng.

\section{Các Phát hiện Chính}

\textbf{Sự đảo ngược accuracy--AUC:} CNN tạo quyết định cứng sắc nét hơn, Swin-Tiny tạo phân phối xác suất được hiệu chỉnh tốt hơn. Mô hình phù hợp phụ thuộc vào nhiệm vụ lâm sàng: cờ nhị phân cần F1 cao (EfficientNet-B0), danh sách xếp hạng rủi ro cần AUC cao (Swin-Tiny).

\textbf{Label Smoothing:} Loại bỏ đồng thời tăng AUC (+0,015) và giảm F1 ($-$3,1\%). Đánh đổi giữa độ sắc nét hiệu chỉnh và chất lượng ranh giới quyết định.

\textbf{Melanoma vẫn khó nhất:} F1 ensemble chỉ 0,744 --- phản ánh độ khó lâm sàng thực sự trong phân tách ung thư hắc tố sớm khỏi nốt ruồi lành tính.

\section{Hạn chế}

\begin{itemize}
    \item \textbf{Thiên kiến nhân khẩu học của tập dữ liệu.} HAM10000 thu thập chủ yếu ở Áo và Mỹ với dân số đa số Fitzpatrick I--III. Triển khai trên bệnh nhân với tông da đậm hơn (Fitzpatrick IV--VI) nhiều khả năng kém hiệu năng, nhất quán với báo cáo thiên kiến tông da ở các hệ thống dermoscopy thương mại. Khóa luận không bao gồm dữ liệu Việt Nam hoặc dữ liệu địa phương đại diện, nên khoảng cách nhân khẩu học được thừa nhận nhưng chưa được đo lường.
    \item \textbf{Khoảng cách modality thu thập.} Đánh giá ngoài phân phối trên PAD-UFES-20 định lượng mức giảm 56 điểm phần trăm accuracy khi cùng một ensemble chỉ-ảnh chạy trên ảnh smartphone lâm sàng thay vì dermoscopy ánh sáng phân cực. Pipeline chưa triển khai được trên ảnh consumer-grade nếu không có domain adaptation, fine-tune in-domain, hoặc tiền xử lý đặc trị cho modality.
    \item \textbf{Phạm vi closed-set bảy lớp.} Classifier bị buộc gán mọi input vào một trong bảy lớp ISIC. Thực hành lâm sàng có hàng chục loại tổn thương, và bất kỳ input nào ngoài bảy lớp huấn luyện (mụn, sẹo, vùng không phải tổn thương) sẽ bị gán nhầm vào lớp gần nhất. Prototype Streamlit giảm thiểu bằng cảnh báo ngưỡng tin cậy, nhưng một bộ phát hiện open-set hoặc OOD chuyên dụng chưa được triển khai.
    \item \textbf{Độ phân giải cố định $224 \times 224$.} Độ phân giải cao hơn, mà các kiến trúc như EfficientNetV2~\cite{tan2021efficientnetv2} thiết kế để khai thác, có thể khôi phục các đặc trưng phân biệt tinh tế bị mất ở kích thước chuẩn, đặc biệt cho các pattern sắc tố và kết cấu tinh tế trong cụm nhầm lẫn MEL/NV/BKL.
    \item \textbf{Tập kiểm tra lớp hiếm nhỏ.} Test set chỉ có 17 Dermatofibroma và 22 Vascular Lesion. Một phân loại sai bổ sung dịch chuyển F1 các lớp này khoảng 0,06, nên các con số per-class cho lớp hiếm mang phương sai cao và không nên diễn giải quá mức.
    \item \textbf{Kiến trúc Transformer đơn.} Chỉ Swin-Tiny được đánh giá. Các biến thể khác (DeiT~\cite{touvron2021deit}, ConvNeXt~\cite{liu2022convnext}, hoặc Swin lớn hơn) có thể cho trade-off accuracy--AUC khác và đặc tính ensemble khác khi kết hợp với các CNN.
    \item \textbf{Không có xác thực lâm sàng tiến cứu.} Các mô hình chưa được triển khai trong workflow patient-facing hoặc đánh giá bởi bác sĩ da liễu thực hành theo thời gian thực. Các nghiên cứu so sánh human--AI và đánh giá chấp nhận của bác sĩ vẫn là điều kiện tiên quyết cho mọi cân nhắc triển khai.
\end{itemize}

\section{Hướng Phát triển}

\begin{itemize}
    \item \textbf{Domain adaptation cho ảnh lâm sàng.} Mức giảm 56 điểm phần trăm OOD trên PAD-UFES-20 thiết lập một mục tiêu định lượng. Style-transfer preprocessing, adversarial domain adaptation, hoặc fine-tune trên một bộ ảnh smartphone gắn nhãn nhỏ đều được kỳ vọng thu hẹp phần đáng kể khoảng cách đó.
    \item \textbf{Open-set và OOD detection.} Bổ sung một component chuyên dụng (Mahalanobis distance trên features backbone, Energy-based score, hoặc tín hiệu tái cấu trúc autoencoder) sẽ cho phép hệ thống từ chối input không phải tổn thương da, thay vì buộc gán bảy lớp mà người dùng phải tự loại trừ.
    \item \textbf{Augmentation theo tông da và phân tích fairness.} Augmentation hue/saturation theo Fitzpatrick, kết hợp với loss term fairness theo subgroup, sẽ trực tiếp đối phó thiên kiến nhân khẩu học thừa hưởng từ phân phối huấn luyện HAM10000. Khi có dữ liệu Việt Nam hoặc dữ liệu địa phương đại diện, các con số fairness có thể được đo lường.
    \item \textbf{Huấn luyện độ phân giải cao hơn.} Các kiến trúc thiết kế cho input lớn hơn có thể khôi phục features bị mất ở $224 \times 224$, đặc biệt cho cụm nhầm lẫn MEL/NV/BKL nơi các khác biệt texture nhỏ mang thông tin chẩn đoán.
    \item \textbf{Huấn luyện trước tự giám sát.} Các phương pháp như DINO~\cite{caron2021dino} có thể trích xuất biểu diễn thị giác mạnh từ các bộ sưu tập dermoscopy không nhãn lớn trước khi fine-tune trên dữ liệu có nhãn, có khả năng giảm yêu cầu dữ liệu hiệu quả và cải thiện tổng quát hóa lớp thiểu số.
    \item \textbf{Phân loại hỗ trợ phân đoạn.} Dự đoán đồng thời ranh giới tổn thương~\cite{hasan2020dsnet} cùng với nhãn bệnh sẽ cho mô hình giám sát không gian rõ ràng cho các features ranh giới liên quan nhất tới phân biệt melanoma--nevus.
    \item \textbf{Knowledge distillation cho triển khai.} Nén ensemble bốn mô hình thành một mô hình nhẹ duy nhất sẽ giữ phần lớn lợi thế accuracy đồng thời giảm chi phí inference về một forward pass, yêu cầu thực tiễn cho triển khai mobile hoặc point-of-care.
    \item \textbf{Đánh giá lâm sàng tiến cứu.} Một nghiên cứu có kiểm soát so sánh output của hệ thống với một hội đồng bác sĩ da liễu thực hành trên tập case đại diện là bước tiếp theo cần thiết trước khi đưa ra bất kỳ tuyên bố triển khai nào.
\end{itemize}
